% Generated draft: requires mathematical and source review.
\documentclass[11pt]{article}
\usepackage[margin=25mm]{geometry}
\usepackage{fontspec}
\setmainfont{Latin Modern Roman}
\usepackage{amsmath,amssymb,mathtools,unicode-math}
\setmathfont{Latin Modern Math}
\usepackage{enumitem,needspace,xurl,hyperref}
\setlength{\parindent}{0pt}
\setlength{\parskip}{5pt}
\setlength{\emergencystretch}{3em}
\newcommand{\N}{\mathbb{N}}
\newcommand{\Z}{\mathbb{Z}}
\newcommand{\Q}{\mathbb{Q}}
\newcommand{\R}{\mathbb{R}}
\newcommand{\C}{\mathbb{C}}
\newcommand{\F}{\mathbb{F}}
\newcommand{\A}{\mathbb{A}}
\newcommand{\PP}{\mathbb{P}}
\newcommand{\E}{\mathbb{E}}
\newcommand{\eps}{\varepsilon}
\newcommand{\dd}{\,\mathrm{d}}
\begin{document}
\section*{UnsolvedMath ID: 30004407}
\textbf{Source problem code:} catalogue entry
\subsection{Definitions and mathematical statement}
Let \(I=[0,1]\), and let \(K(I^3)\) denote the knot-isotopy classes intended by the supplied source. Write \(U\) for the unknot and \(K\mathbin{\#}L\) for connected sum. The notation \(T_n(I,I^3)\) refers to the intended \(n\)-th stage of the embedding-calculus tower, and \(\operatorname{ev}_n\) denotes its evaluation map. The record does not specify the embedding-space model, regularity, boundary conditions, or precise tower construction. In particular, the notation suggests knots represented by embedded intervals in the cube, but the required endpoint and boundary-collar conditions are not supplied. These conventions must be recovered from the cited report before the statement is fully specified.

For an abelian group \(A\), an additive knot invariant is a function \(v:K(I^3)\to A\) satisfying
\[
v(U)=0,\qquad v(K\mathbin{\#}L)=v(K)+v(L).
\]
Such an invariant has Vassiliev type at most \(m\) if its alternating extension vanishes on every singular knot with \(m+1\) transverse double points, in the same intended knot category. Explicitly, for every such singular knot \(S\), with its two local resolutions at each double point labeled \(+\) and \(-\), this means
\[
\sum_{\varepsilon\in\{+,-\}^{m+1}}
(-1)^{\#\{i:\varepsilon_i=-\}}v(S_\varepsilon)=0.
\]

\textbf{Question.} For every integer \(n\geq1\), is
\[
q_n:=\pi_0\operatorname{ev}_n:
K(I^3)\longrightarrow G_n,
\qquad G_n:=\pi_0T_n(I,I^3),
\]
a universal additive Vassiliev invariant of type at most \(n-1\) over \(\mathbb Z\)? With the intended abelian-group structure on \(G_n\), the universal-property formulation asks that \(q_n\) itself be additive and of type at most \(n-1\), and that for every abelian group \(A\) and every additive invariant \(v:K(I^3)\to A\) of type at most \(n-1\), there exist a unique group homomorphism \(\overline v:G_n\to A\) such that
\[
v=\overline v\circ q_n.
\]
The group structure and its compatibility with evaluation are required for this formulation, but are not defined in the supplied record.

The source also presents the following as an equivalent formulation: \(q_n\) is surjective and
\[
\{K\in K(I^3):q_n(K)=q_n(U)\}
=
\{K\in K(I^3):K\sim_n U\},
\]
where \(\sim_n\) is the source's relation of \(n\)-equivalence. Under the intended additive structure, \(q_n(U)=0\), so the left-hand side is the stated kernel. The record does not define \(n\)-equivalence; its convention and the hypotheses justifying the asserted equivalence of formulations remain unresolved.

The supplied dated assessment labels the problem partially solved: it reports that factorization and surjectivity are proved, that the resulting map is an isomorphism rationally for every \(n\) and \(p\)-adically for \(n\leq p+2\), and that the integral isomorphism conjecture remains open. Here \(p\) is understood as a prime in the reported \(p\)-adic qualification. The assessment does not identify the factorization map, the relevant quotient, or the precise rational and \(p\)-adic constructions. These are supplied literature claims, not independently verified assertions in this draft.
\subsection{Short English statement}
For each \(n\geq1\), does the evaluation map to the \(n\)-th embedding-calculus stage universally encode all additive integral finite-type knot invariants of degree at most \(n-1\)? The source equivalently asks for surjectivity and for the unknot fiber to consist exactly of knots \(n\)-equivalent to the unknot. Precise knot, tower, and equivalence conventions require review. Supplied dated curation reports partial results and leaves the integral conjecture open.
\subsection{Sources}
\begin{itemize}
\item Supplied local UnsolvedMath registry record, permanent ID \(30004407\), source problem code \texttt{catalogue entry}, titled \emph{Universal Additive Vassiliev Invariants from Embedding Calculus}, in the Oberwolfach Reports Open Problems collection. No individual proposer is supplied; the proposed year is recorded as 2020.
\item Supplied source citation: \emph{Low-dimensional Topology (2021)}, Oberwolfach Reports, pp.\ 465--516, DOI \texttt{10.4171/owr/2020/8}. \url{https://doi.org/10.4171/owr/2020/8}. The supplied background describes Oberwolfach Reports as documenting mathematical workshops and their open-problem discussions.
\item Supplied literature source labeled \emph{Oberwolfach Report 8/2020}: \url{https://ems.press/content/serial-article-files/46844}.
\item Supplied literature source also labeled \emph{Oberwolfach Report 8/2020}: \url{https://doi.org/10.4171/owr/2020/8}.
\item Literature assessment date: \texttt{2026-08-22}. The partial-solution assessment and its factorization, surjectivity, rational, and \(p\)-adic claims are supplied dated curation derived from the OWR material, not independently verified here. The record explicitly requests independent verification of time-sensitive status and cited claims.
\item The registry labels the statement \texttt{corrected\_verified} and says that mathematical content was preserved while notation, typography, omitted definitions, or duplicate clauses were repaired. This is a supplied editorial qualification and does not establish mathematical correctness or resolve the missing conventions identified above.
\end{itemize}
\end{document}
